1.
Introduction générale
Notion d'état d'un système physique :
-Cas classique et cas quantique
-Évolution d'un système
Notion d'"espace des états":
Principe de combinaison
Structure de l'espace des états : espace vectoriel
Corps des scalaires : nombres complexes
Notation des vecteurs d'états : "kets"
ATTENTION : mise en garde essentielle à propos des combinaisons linéaires d'états

2.
Évolution d'un système physique
Opérateur d'évolution
Comment distinguer les états physiques ?
- États à valeur bien définie d'une grandeu donnée
- États sans valeur bien définie
- Grandeurs physiques et combinaisons linéaires
Problématique générale de la mesure :
-interaction système/instrument
interprétation de Copenhague, interprétation d'Everett

3.
Problématique générale de la mesure (fin)
- Indéterminisme, aléatoire, probabilités
- Règle de Born
- Impossibilité de prendre connaissance d'un état initialement inconnu
Structure herlmitienne de l'espace des états
- produit scalaire sur un espace vectoriel réel
Attention : la dernière demi-heure est sans son, en raison d'une défaillance de l'alimentation de système audio... Toute mes excuses !

4.

5.
Adjoint d'un opérateur
Opérateurs auto-adjoints
Construction de l'espace des états d'un système physique :
- base orthonormée de vecteurs à valeur bien définie d'une grandeur physique
-"coordonnées" d'un vecteur d'état
- relation de fermeture, projecteur orthogonal
- mesure d'une grandeur physique

6.
=== Fin du chapitre 1
Espace des états d'un système quantique
Grandeurs physiques complémentaires
Notion d'observable physique
"Compatibilité" entre observables

=== Début du chapitre 2
Systèmes quantiques élémentaires
Notion de qubit

7.
Systèmes quantiques élémentaires : qubit
- représentation d'un qubit sur la sphère de Bloch
- représentation d'un ket en vecteur colonne
- représentation matricielle d'un opérateur
- observable sur un qubit (opérateur auto-adjoint)
- matrices de Pauli, valeurs et vecteurs propres
Expérience de Stern et Gerlach (1922)
- Rappel : moment magnétique
- Dispositif de Stern et Gerlach, analyse classique

8.
Expérience de Stern et Gerlach (1922)
